\subsection{\texorpdfstring{类型为 \(\mathbf{B}_l\) 的代数 ( \(l \geq 1\) )}{类型为 \textbackslash mathbf\{B\}\_l 的代数 ( l \textbackslash geq 1 )}}

\begin{enumerate}
\def\labelenumi{(\Roman{enumi})}
\tightlist
\item
  设 V 为有限维向量空间，\(\Psi\) 为 V 上的非退化对称双线性型。所有满足 \(\Psi(xv,v') + \Psi(v, xv') = 0\)（对所有 \(v, v' \in V\)）的 V 的自同态 x 构成 \(\mathfrak{sl}(V)\) 的李子代数；当 dim \(V \neq 2\) 时，它是半单的（第 I 章，§6，第7节，命题9）。我们将其记为 \(\mathfrak{o}(\Psi)\)，称其为与 \(\Psi\) 相关的正交李代数。
\end{enumerate}

设 V 的维数为奇数 \(2l + 1 \geq 3\)，且 \(\Psi\) 的指标为最大指标 l。记 Q 为与 \(\Psi\) 相关的二次型。有 \(\mathrm{Q}(x)=\frac{1}{2}\varPsi(x,x)\) 对 \(x\in V\) 成立。由《代数》第 IX 章，§4，第2节，V 可写成两个极大全迷向子空间 F 和 \(F'\) 与 \(F+F'\) 的正交补 G 的直和；G 非迷向且为一维。将 \(\varPsi\) 乘以一个非零常数后，不妨设存在 \(e_{0}\in G\)，使得 \(\varPsi(e_{0},e_{0})=-2\)。另一方面，F 和 \(F'\) 通过 \(\varPsi\) 处于对偶关系；令 \((e_{i})_{1\leq i\leq l}\) 为 F 的一个基，令 \((e_{-i})_{1\leq i\leq l}\) 为 \(F'\) 的对偶基。于是

\[
(e _ {1}, \ldots , e _ {l}, e _ {0}, e _ {- l}, \ldots , e _ {- 1})
\]

是 V 的一个基；有

\[
\mathrm{Q} \left(\sum x _ {i} e _ {i}\right) = - x _ {0} ^ {2} + \sum_ {i = 1} ^ {i = l} x _ {i} x _ {- i}
\]

并且相对于此基，\(\varPsi\) 的矩阵是阶为 \(2l + 1\) 的方阵

\[
S = \left( \begin{array}{c c c} 0 & 0 & s \\ 0 & - 2 & 0 \\ s & 0 & 0 \end{array} \right), \qquad s = \left( \begin{array}{c c c c c} 0 & 0 & \dots & 0 & 1 \\ 0 & 0 & \dots & 1 & 0 \\ \vdots & \vdots & \ddots & \vdots & \vdots \\ 0 & 1 & \dots & 0 & 0 \\ 1 & 0 & \dots & 0 & 0 \end{array} \right),
\]

其中 s 是阶为 l 的方阵，除第二对角线 \(^{6}\) 上的元素等于 1 外，其余元素全为 0。具有上述性质的 V 的基称为 Witt 基。于是，代数 \(\mathfrak{g} = \mathfrak{o}(\varPsi)\) 可与方阵代数 \(\mathfrak{o}_{S}(2l + 1, k)\) 识别，后者由阶为 \(2l + 1\) 的方阵 a 组成，且满足 \(a = -S^{-1}aS\)（《代数》第 IX 章，§1，第10节，公式（50））。直接计算表明，g 是如下形式的矩阵的集合：

\[
\left( \begin{array}{c c c} A & 2 s ^ {t} x & B \\ y & 0 & x \\ C & 2 s ^ {t} y & D \end{array} \right)\tag{2}
\]

其中 x 和 y 是具有 1 行 l 列的矩阵，A、B、C、D 是阶为 l 的方阵，并满足 \(B = -s^{t}Bs\)、\(C = -s^{t}Cs\) 和 \(D = -s^{t}As\)。由于从 \(A \mapsto s^{t}As\) 到其自身的映射 \(\mathbf{M}_{l}(k)\) 是关于第二对角线的对称，故

\[
\dim \mathfrak {g} = 2 l + l ^ {2} + 2 \frac {l (l - 1)}{2} = l (2 l + 1).
\]

令 \(\mathfrak{h}\) 为 \(\mathfrak{g}\) 的对角元素的集合。这是 \(\mathfrak{g}\) 的交换子代数，其基由以下元素组成：

\footnotetext[6]{方阵 \((a_{ij})_{1\leq i,j\leq n}\) 的第二对角线是满足 \(a_{ij}\) 的 \(i+j=n+1\) 所成的族。}

\[
H _ {i} = E _ {i, i} - E _ {- i, - i} \quad (1 \leq i \leq l).
\]

令 \((\varepsilon_{i})\) 为与 \(h^{*}\) 对偶的 \((H_{i})\) 的基。置

\[
\left\{ \begin{array}{l l l} X _ {\varepsilon_ {i}} & = 2 E _ {i, 0} + E _ {0, - i} & (1 \leq i \leq l) \\ X _ {- \varepsilon_ {i}} & = - 2 E _ {- i, 0} - E _ {0, i} & (1 \leq i \leq l) \\ X _ {\varepsilon_ {i} - \varepsilon_ {j}} & = E _ {i, j} - E _ {- j, - i} & (1 \leq i <   j \leq l) \\ X _ {\varepsilon_ {j} - \varepsilon_ {i}} & = - E _ {j, i} + E _ {- i, - j} & (1 \leq i <   j \leq l) \\ X _ {\varepsilon_ {i} + \varepsilon_ {j}} & = E _ {i, - j} - E _ {j, - i} & (1 \leq i <   j \leq l) \\ X _ {- \varepsilon_ {i} - \varepsilon_ {j}} & = - E _ {- j, i} + E _ {- i, j} & (1 \leq i <   j \leq l). \end{array} \right.\tag{3}
\]

容易验证，这些元素构成 \(\mathfrak{h}\) 中 \(\mathfrak{g}\) 的一个补空间的基，并且对于 \(h\in \mathfrak{h}\)，有

\[
[ h, X _ {\alpha} ] = \alpha (h) X _ {\alpha}\tag{4}
\]

对所有 \(\alpha \in \mathbb{R}\) 都成立，其中 \(\mathbb{R}\) 是由 \(\pm \varepsilon_i\) 以及 \(\pm \varepsilon_i \pm \varepsilon_j\)（\(1 \leq i < j \leq l\)）组成的集合。因此，\(\mathfrak{h}\) 等于其在 \(\mathfrak{g}\) 中的正规化子，因而是 \(\mathfrak{g}\) 的嘉当子代数；\(\mathfrak{h}\) 是分裂的，并且 \((\mathfrak{g}, \mathfrak{h})\) 的根正是 \(\mathbb{R}\) 的元素。\(\mathbb{R}\) 的根系 \((\mathfrak{g}, \mathfrak{h})\) 在 \(\mathrm{B}_l\) 时为 \(l \geq 2\) 型，在 \(\mathrm{A}_1\) 时为 \(\mathrm{B}_1\) 型（也称 \(l = 1\) 型）（第 VI 章，§4，第5.I节，推广到 \(l = 1\) 的情形）。因此，\(\mathfrak{g}\) 是 \(\mathrm{B}_l\) 型的可分裂单李代数。

正交代数 \(\mathfrak{o}(\varPsi)\) 的每个分裂嘉当子代数都是 h 在 \(\mathfrak{o}(\varPsi)\) 的某个初等自同构下的变换，因而也是在 \(\mathbf{O}(\varPsi)\) 的某个元素下的变换（参见~(VII)）；所以它就是集合 \(h_{\beta}\)，其中的元素 g 相对于 V 的某个 Witt 基 \(\beta\) 的矩阵为对角矩阵。立即验证可知，在 \(h_{\beta}\) 下不变的唯一向量子空间是由 \(\beta\) 的某个子集生成的子空间。

若 \(l = 1\)，代数 \(\mathfrak{o}(\varPsi)\) 和 \(\mathfrak{sl}(2, k)\) 具有相同的根系，因而同构（另参见~§1，习题16）。下文设 \(l \geq 2\)。

\begin{enumerate}
\def\labelenumi{(\Roman{enumi})}
\setcounter{enumi}{1}
\tightlist
\item
  根系 \(R^{\vee}\) 由第 VI 章，§4，第5.V节确定，我们得到
\end{enumerate}

\[
H _ {\varepsilon_ {i}} = 2 H _ {i}, \quad H _ {\varepsilon_ {i} - \varepsilon_ {j}} = H _ {i} - H _ {j}, \quad H _ {\varepsilon_ {i} + \varepsilon_ {j}} = H _ {i} + H _ {j}.
\]

\begin{enumerate}
\def\labelenumi{(\Roman{enumi})}
\setcounter{enumi}{2}
\tightlist
\item
  置 \(\alpha_{1} = \varepsilon_{1} - \varepsilon_{2},\ldots ,\alpha_{l - 1} = \varepsilon_{l - 1} - \varepsilon_{l},\alpha_{l} = \varepsilon_{l}\)。由第 VI 章，§4，第5.II节，\((\alpha_{1},\dots ,\alpha_{l})\) 是 R 的一个基 B；相对于 B 的正根是 \(\varepsilon_{i}\) 以及 \(\varepsilon_{i}\pm \varepsilon_{j}\)（\(i < j\)）。相应的 Borel 子代数 \(\mathfrak{b}\) 是 \(\mathfrak{g}\) 中上三角矩阵的集合。
\end{enumerate}

立即验证可知，在 b 下稳定且不同于 \(\{0\}\) 和 V 的唯一向量子空间，是与基 \((e_{i})\) 对应的极大旗中的元素，即全迷向子空间 \(V_{1},\ldots,V_{l}\)，其中 \(V_{i}\) 由 \(e_{1},\ldots,e_{i}\) 生成，以及它们的正交补 \(V_{-1},\ldots,V_{-i}\)：\(V_{-i}\) 的正交补 \(V_{i}\) 由 \(e_{1},\ldots,e_{l},e_{0},e_{-l},\ldots,e_{-i-1}\) 生成，且不是全迷向的。另一方面，如果 g 的一个元素使某个向量子空间稳定，它也使其正交补稳定。因此，\(\mathfrak{b}\) 是使旗 \(\mathfrak{g}\) 的各元素稳定的 \(\{\mathrm{V}_1,\dots ,\mathrm{V}_l\}\) 中元素的集合。

如果旗的每个元素都是全迷向的，就称该旗为迷向旗。旗 \(\{V_{1},\ldots,V_{l}\}\) 是极大迷向旗。由于群 \(\mathbf{O}(\varPsi)\) 既在 g 的 Borel 子代数上（参见~(VII)）传递作用，也在极大迷向旗上（《代数》第 IX 章，§4，第3节，定理1）传递作用，我们看到，对于 V 中任意极大迷向旗 \(\delta\)，使 \(b_{\delta}\) 的各元素稳定的 g 中元素的集合 \(\delta\) 是 g 的一个 Borel 子代数，并且映射 \(\delta\mapsto b_{\delta}\) 是从极大迷向旗集合到 g 的 Borel 子代数集合的双射。

令 \(\delta\) 为迷向旗，令 \(\mathfrak{p}_{\delta}\) 为使 \(\mathfrak{g}\) 的各元素稳定的 \(\delta\) 中元素的集合。若 \(\delta \subset \{\mathrm{V}_1, \ldots, \mathrm{V}_l\}\)，则 \(\mathfrak{p}_{\delta}\) 是包含 \(\mathfrak{g}\) 的 \(\mathfrak{b}\) 的一个抛物子代数，并且容易验证，在 \(\neq \{0\}\) 下稳定且不等于 \(\mathfrak{p}_{\delta}\) 的唯一全迷向子空间是 \(\delta\) 的元素。这给出包含 \(2^l\) 的 \(\mathfrak{g}\) 的 \(\mathfrak{b}\) 个抛物子代数。如上所述，映射 \(\delta \mapsto \mathfrak{p}_{\delta}\) 是从 V 中迷向旗的集合到 \(\mathfrak{g}\) 的抛物子代数集合的双射。此外，当且仅当 \(\mathfrak{p}_{\delta} \subset \mathfrak{p}_{\delta'}\) 时，\(\delta \supset \delta'\)。

\begin{enumerate}
\def\labelenumi{(\Roman{enumi})}
\setcounter{enumi}{3}
\tightlist
\item
  与 \(\alpha_{1},\ldots,\alpha_{l}\) 对应的基本权由第 VI 章，§4，第5.VI节给出：
\end{enumerate}

\[
\begin{array}{l} \varpi_ {i} = \varepsilon_ {1} + \dots + \varepsilon_ {i} (1 \leq i \leq l - 1) \\ \varpi_ {l} = \frac {1}{2} (\varepsilon_ {1} + \dots + \varepsilon_ {l}). \end{array}
\]

令 \(\sigma\) 为 \(\mathfrak{g}\) 在 V 上的恒等表示。外幂 \(\bigwedge^{r}\sigma\) 作用于 \(E = \bigwedge^{r}V\)。若 \(h\in \mathfrak{h}\)，

\[
\begin{array}{l} \sigma (h). e _ {i} = \varepsilon_ {i} (h) e _ {i} \quad \text { for } 1 \leq i \leq l \\ \sigma (h). e _ {0} = 0 \\ \sigma (h). e _ {- i} = - \varepsilon_ {i} (h) e _ {- i} \quad \text { for } 1 \leq i \leq l. \end{array}
\]

由此，对于 \(1 \leq r \leq l\)，\(\varepsilon_{1} + \cdots + \varepsilon_{r}\) 是 \(\bigwedge^{r}\sigma\) 的最高权，权为 \(\varepsilon_{1} + \cdots + \varepsilon_{r}\) 的元素正是与 \(e_{1} \wedge \cdots \wedge e_{r}\) 成比例的元素。我们将证明，当 \(1 \leq r \leq l - 1\) 时，表示 \(\bigwedge^{r}\sigma\) 是最高权为 \(\varpi_{r}\) 的基本表示。为此，只需证明 \(\bigwedge^{r}\sigma\) 在 \(0 \leq r \leq 2l + 1\) 时不可约。但是，在 \(\Phi\) 上定义的双线性型 \(\bigwedge^{r}V \times \bigwedge^{2l+1-r}V\) 为

\[
x \wedge y = \varPhi (x, y) e _ {1} \wedge \dots \wedge e _ {l} \wedge e _ {0} \wedge e _ {- l} \wedge \dots \wedge e _ {- 1}
\]

它在 g 下不变，并使 \(\bigwedge^{r}V\) 与 \(\bigwedge^{2l+1-r}V\) 处于对偶关系。因此，表示 \(\bigwedge^{2l+1-r}\sigma\) 是 \(\bigwedge^{r}\sigma\) 的对偶，只需证明 \(\bigwedge^{r}\sigma\) 在 \(0\leq r\leq l\) 时不可约；等价地，只需证明 \(T_{r}\) 中包含 \(\bigwedge^{r}V\) 且在 g 下稳定的最小子空间 \(e_{1}\wedge\cdots\wedge e_{r}\) 就是整个 \(\bigwedge^{r}V\)。当 r=0 和 r=1 时这立即成立（参见~公式（2））。当 r=2（从而 \(l\geq2\)）时，\(\bigwedge^{2}\sigma\) 与 g 的伴随表示（后者不可约）具有相同的维数 \(l(2l + 1)\) 和相同的最高权 \(\varepsilon_1 + \varepsilon_2\)（第 VI 章，§4，第5.IV节）。因此，\(\bigwedge^2\sigma\) 等价于伴随表示，因而不可约。这就证明了 \(l = 1\) 和 \(l = 2\) 时的断言。

现在对 \(l\) 作归纳，并设 \(l \geq r \geq 3\)。先注意，若 W 是 V 的一个奇数维非迷向子空间，正交补为 \(\mathrm{W}'\)，则 \(\Psi_{\mathrm{W}}\)（即 \(\Psi\) 在 W 上的限制）非退化，并且 \(\mathfrak{o}(\Psi_{\mathrm{W}})\) 可与 \(\mathfrak{g}\) 中由在 \(\mathrm{W}'\) 上为零的元素组成的子代数识别。若 \(\dim \mathrm{W} < \dim \mathrm{V}\)，且 \(\Psi_{\mathrm{W}}\) 的指标为最大指标，则归纳假设表明：如果 \(\mathrm{T}_r\) 含有形如 \(w' \wedge w\) 的非零元素，其中 \(w' \in \bigwedge^{r - k} \mathrm{W}'\)，\(w \in \bigwedge^k \mathrm{W}\)（\(0 \leq k \leq r\)），那么 \(\mathrm{T}_r\) 含有 \(w' \wedge \bigwedge^k \mathrm{W}\)：事实上，有 \(a.(w' \wedge w) = w' \wedge a.w\)，对所有 \(a \in \mathfrak{o}(\Psi_{\mathrm{W}})\)。我们对 \(p \in [0, r]\) 作归纳，证明 \(\mathrm{T}_r\) 含有元素

\[
x = e _ {i _ {1}} \wedge \dots \wedge e _ {i _ {r - p}} \wedge e _ {j _ {1}} \wedge \dots \wedge e _ {j _ {p}}
\]

当 \(1 \leq i_{1} < \cdots < i_{r-p} \leq l\) 且 \(-l \leq j_{1} < \cdots < j_{p} \leq 0\) 时。对于 p = 0，这源于 \(\mathfrak{gl}(F)\) 在 \(\bigwedge^{r} F\) 上作用的不可约性（第1节），因为 g 包含使 \(F = V_{l} = \sum_{i=1}^{l} ke_{i}\) 不变并在其上诱导任意自同态的元素（参见~公式（2））。若 p = 1，取 \(q \in (1, l)\)，使得 \(q \neq -j_{1}\)，且存在 \(\lambda \in [1, r - p]\) 使 \(q = i_{\lambda}\)；若 \(p \geq 2\)，取 \(q \in [1, l]\)，使得 \(-q \in \{j_{1}, \ldots, j_{p}\}\)。必要时置换 \(e_{i}\)，可设 q = 1。现在取 W 为 \(W' = ke_{1} + ke_{-1}\) 的正交补。若 p = 1，则有 \(x \in e_{1} \wedge \bigwedge^{r-1} W\)；由于 \(T_{r}\) 含有 \(e_{1} \wedge \cdots \wedge e_{r}\)，我们看到 \(T_{r}\) 含有 x。若 \(p \geq 2\)，则或者 \(x \in e_{-1} \wedge \bigwedge^{r-1} W\)，或者 \(x \in e_{1} \wedge e_{-1} \wedge \bigwedge^{r-2} W\)；由于根据归纳假设，\(T_{r}\) 含有 \(e_{-1} \wedge e_{2} \wedge \cdots \wedge e_{r-1}\) 和 \(e_{-1} \wedge e_{1} \wedge e_{2} \wedge \cdots \wedge e_{r-2}\)，我们看到 \(T_{r}\) 含有 x，从而证明完毕。

关于 \(\bigwedge^{r}\sigma\) 不可约性的另一证明，见习题6。

现在确定最高权为 \(\varpi_{l}\) 的基本表示。

引理 1. 设 \(\mathrm{V}\) 为有限维向量空间，\(\mathbb{Q}\) 为 \(\mathrm{V}\) 上的非退化二次型，\(\Psi\) 为与 \(\mathbb{Q}\) 相关的对称双线性型，\(\mathrm{C(Q)}\) 为相对于 \(\mathrm{V}\) 的 \(\mathbb{Q}\) 的 Clifford 代数，\(f_0\) 为以下典则映射的复合

\[
\mathfrak {o} (\Psi) \longrightarrow \mathfrak {g l} (V) \longrightarrow V \otimes V ^ {*} \longrightarrow V \otimes V \longrightarrow C ^ {+} (Q)
\]

（第一个是典则嵌入，第三个由对应于 \(\mathrm{V}^*\) 的从 \(\mathrm{V}\) 到 \(\Psi\) 的典则同构定义，第四个由 \(\mathrm{C(Q)}\) 中的乘法定义，参见~《代数》第 IX 章，§ 9，第1节）。置 \(f = \frac{1}{2} f_0\)。

\begin{enumerate}
\def\labelenumi{(\roman{enumi})}
\item
  若 \((e_r), (e_r')\) 是 V 的满足 \(\Psi(e_r, e_s') = \delta_{rs}\) 的基，则对所有 \(f_0(a) = \sum_r (ae_r)e_r'\)，有 \(a \in \mathfrak{o}(\Psi)\)。
\item
  若 \(a, b \in \mathfrak{o}(\Psi)\)，则有 \(\sum_{r}(ae_{r})(be_{r}^{\prime}) = -\sum_{r}(abe_{r})e_{r}^{\prime}\)。
\item
  若 \(a \in \mathfrak{o}(\varPsi)\) 且 \(v \in V\)，则有 \([f(a), v] = av\)。
\item
  若 \(a, b \in \mathfrak{o}(\varPsi)\)，则有 \([f(a), f(b)] = f([a, b])\)。
\item
  \(f(\mathfrak{o}(\Psi))\) 生成结合代数 \(\mathrm{C}^{+}(\mathrm{Q})\)。
\item
  设 \(\mathrm{N}\) 为左 \(\mathrm{C}^{+}(\mathrm{Q})\)-模，\(\rho\) 为从 \(\mathrm{C}^{+}(\mathrm{Q})\) 到 \(\operatorname{End}_k(\mathrm{N})\) 的相应同态。则 \(\rho \circ f\) 是 \(\mathfrak{o}(\Psi)\) 在 \(\mathrm{N}\) 上的表示。若 \(\mathrm{N}\) 是单模，则 \(\rho \circ f\) 不可约。
\end{enumerate}

断言 (i) 是显然的。若 \(a, b \in \mathfrak{o}(\varPsi)\)，记 \(\varPsi(x,y) = \langle x,y\rangle\)，则有：

\[
\begin{array}{c} \sum_ {r} (a e _ {r}) (b e _ {r} ^ {\prime}) = \sum_ {r, s, t} \langle a e _ {r}, e _ {s} ^ {\prime} \rangle \langle b e _ {r} ^ {\prime}, e _ {t} \rangle e _ {s} e _ {t} ^ {\prime} = \sum_ {r, s, t} \langle e _ {r}, a e _ {s} ^ {\prime} \rangle \langle e _ {r} ^ {\prime}, b e _ {t} \rangle e _ {s} e _ {t} ^ {\prime} \\ = \sum_ {s, t} \langle a e _ {s} ^ {\prime}, b e _ {t} \rangle e _ {s} e _ {t} ^ {\prime} = - \sum_ {s, t} \langle e _ {s} ^ {\prime}, a b e _ {t} \rangle e _ {s} e _ {t} ^ {\prime} = - \sum_ {t} (a b e _ {t}) e _ {t} ^ {\prime} \end{array}
\]

这证明了 (ii)。接着，对所有 \(v \in V\)，由 (i) 得

\[
\begin{array}{l} [ f (a), v ] = \frac {1}{2} \sum_ {r} ((a e _ {r}) e _ {r} ^ {\prime} v - v (a e _ {r}) e _ {r} ^ {\prime}) \\ \qquad = \frac {1}{2} \sum_ {r} ((a e _ {r}) e _ {r} ^ {\prime} v + (a e _ {r}) v e _ {r} ^ {\prime} - (a e _ {r}) v e _ {r} ^ {\prime} - v (a e _ {r}) e _ {r} ^ {\prime}) \\ \qquad = \frac {1}{2} \sum_ {r} ((a e _ {r}) \langle e _ {r} ^ {\prime}, v \rangle - \langle a e _ {r}, v \rangle e _ {r} ^ {\prime}) \\ \qquad = \frac {1}{2} a \left(\sum_ {r} \langle e _ {r} ^ {\prime}, v \rangle e _ {r}\right) + \frac {1}{2} \sum_ {r} \langle e _ {r}, a v \rangle e _ {r} ^ {\prime} = \frac {1}{2} a v + \frac {1}{2} a v = a v, \end{array}
\]

这证明了 (iii)。于是

\[
\begin{array}{l} [ f (a), f (b) ] = \left[ f (a), \frac {1}{2} \sum_ {r} (b e _ {r}) e _ {r} ^ {\prime} \right] \quad \text { by   (i) } \\ \quad = \frac {1}{2} \sum_ {r} ([ f (a), b e _ {r} ] e _ {r} ^ {\prime} + (b e _ {r}) [ f (a), e _ {r} ^ {\prime} ]) \\ \quad = \frac {1}{2} \sum_ {r} ((a b e _ {r}) e _ {r} ^ {\prime} + (b e _ {r}) (a e _ {r} ^ {\prime})) \quad \text { by   (iii) } \\ \quad = \frac {1}{2} \sum_ {r} ((a b e _ {r}) e _ {r} ^ {\prime} - (b a e _ {r}) e _ {r} ^ {\prime}) \quad \text { by   (ii) } \\ \quad = f ([ a, b ]) \quad \text { by   (i) } \end{array}
\]

这证明了 (iv)。为证明 (v)，通过扩张标量，可设 \(k\) 代数闭。于是取 V 的一个基 \((e_r)\)，使 \(\varPsi(e_r, e_s) = \delta_{rs}\)，从而 \(e_r' = e_r\)。若 \(i \neq j\)，则 \(\mathrm{E}_{ij} - \mathrm{E}_{ji} \in \mathfrak{o}(\varPsi)\)，且

\[
f (\mathrm{E} _ {i j} - \mathrm{E} _ {j i}) = \frac {1}{2} (e _ {i} e _ {j} - e _ {j} e _ {i}) = e _ {i} e _ {j};
\]

但 \(e_{i}e_{j}\) 生成 \(\mathrm{C}^{+}(\mathrm{Q})\)。

断言 (vi) 由 (iv) 和 (v) 得出。

Q.E.D.

现在回顾相关记号。令 \(\tilde{V} = F + F'\)、\(\tilde{Q}\)（分别为 \(\tilde{\Psi}\)）表示 Q（分别为 \(\Psi\)）在 \(\tilde{V}\) 上的限制。于是，\(\tilde{Q}\) 是定义在 \(\tilde{V}\) 上、维数为 2l 且指标为最大指标 l 的非退化二次型，而 Clifford 代数 \(C(\tilde{Q})\) 是维数为 \(2^{2l}\) 的中心单代数。令 N 为极大迷向子空间 \(F'\) 的外代数，它由 \(e_{-1}, \ldots, e_{-l}\) 生成。利用 \(F'\) 将 F 识别为其对偶空间，并借助 \(\Psi\)，对于 \(x \in F'\)（分别对于 \(y \in F\)），在 N 中将以 x（分别以 y）作左外积（分别作左内积）记为 \(\lambda(x)\)（分别记为 \(\lambda(y)\)）；若 \(a_{1}, \ldots, a_{k} \in F'\)，则

\[
\lambda (x). (a _ {1} \wedge \dots \wedge a _ {k}) = x \wedge a _ {1} \wedge \dots \wedge a _ {k}
\]

\[
\lambda (y). \left(a _ {1} \wedge \dots \wedge a _ {k}\right) = \sum_ {i = 1} ^ {k} (- 1) ^ {i - 1} \Psi \left(a _ {i}, y\right) a _ {1} \wedge \dots \wedge a _ {i - 1} \wedge a _ {i + 1} \wedge \dots \wedge a _ {k}.
\]

容易验证 \(\lambda (x)^{2} = \lambda (y)^{2} = 0\)，并且

\[
\lambda (x) \lambda (y) + \lambda (y) \lambda (x) = \Psi (x, y). 1.
\]

由此（《代数》第 IX 章，§9，第1节），存在唯一的同态（仍记为 \(\lambda\)）从 \(\mathrm{C}(\tilde{\mathbf{Q}})\) 到 \(\mathrm{End}(\mathrm{N})\)，延拓映射 \(\lambda : \mathrm{F} \cup \mathrm{F}' \to \mathrm{End}(\mathrm{N})\)。由于 \(\dim \mathrm{N} = 2^l\)，且 \(\mathrm{C}(\tilde{\mathbf{Q}})\) 只有一类维数为 \(2^l\) 的单模（《代数》第 IX 章，§9，第4节，定理2），由 \(\mathrm{C}(\tilde{\mathbf{Q}})\) 定义的 \(\mathrm{N}\) 在 \(\lambda\) 上的表示是不可约表示，并且是 \(\mathrm{C}(\tilde{\mathbf{Q}})\) 的旋量表示（同上引）。

现在考虑从 \(\mu:v\mapsto e_{0}v\) 到 \(\tilde{\mathrm{V}}\) 的映射 \(\mathrm{C}^{+}(\mathrm{Q})\)。对于 \(v\in\tilde{V}\)，有

\[
(e _ {0} v) ^ {2} = - e _ {0} ^ {2} v ^ {2} = - \mathrm{Q} (e _ {0}) \mathrm{Q} (v) = \mathrm{Q} (v) = \tilde {\mathrm{Q}} (v)
\]

并且 \(\mu\) 唯一延拓为同态（仍记为 \(\mu\)），从 \(\mathrm{C}(\tilde{\mathbb{Q}})\) 到 \(\mathrm{C}^{+}(\mathrm{Q})\)。由于 \(\mathrm{C}(\tilde{\mathbb{Q}})\) 是单代数，且

\[
\dim \mathrm{C} ^ {+} (\mathrm{Q}) = \dim \mathrm{C} (\tilde {\mathrm{Q}}) = 2 ^ {2 l},
\]

可知 \(\mu\) 是同构。因此，\(\lambda \circ \mu^{-1}\) 在 N 上定义了一个单 \(\mathrm{C}^{+}(\mathrm{Q})\)-模结构，而 \(\rho = \lambda \circ \mu^{-1} \circ f\) 是 g 在 N 上的不可约表示（引理1 (vi)）。

另一方面，根据引理1 (i)，有

\[
f (H _ {i}) = \frac {1}{2} (e _ {i} e _ {- i} - e _ {- i} e _ {i}).
\]

由于 \(e_{i}e_{-i} = -e_{0}^{2}e_{i}e_{-i} = e_{0}e_{i}e_{0}e_{-i}\)，且 \(e_{i}e_{-i} + e_{-i}e_{i} = 1\)，有

\[
\mu^ {- 1} \circ f (H _ {i}) = \frac {1}{2} - e _ {- i} e _ {i} = - \frac {1}{2} + e _ {i} e _ {- i}.
\]

由此，对于 \(1 \leq i_{1} < \cdots < i_{k} \leq l\)，有：

\[
\rho (H _ {i}) (e _ {- i _ {1}} \wedge \dots \wedge e _ {- i _ {k}}) = \left\{ \begin{array}{l l} - \frac {1}{2} e _ {- i _ {1}} \wedge \dots \wedge e _ {- i _ {k}} & \text { if } i \in \{i _ {1}, \ldots , i _ {k} \} \\ \frac {1}{2} e _ {- i _ {1}} \wedge \dots \wedge e _ {- i _ {k}} & \text { if } i \notin \{i _ {1}, \ldots , i _ {k} \} \end{array} \right.
\]

并且对于 \(h\in \mathfrak{h}\)

\[
\begin{array}{l} \rho (h) (e _ {- i _ {1}} \wedge \dots \wedge e _ {- i _ {k}}) \\ = (\frac {1}{2} (\varepsilon_ {1} + \dots + \varepsilon_ {l}) - (\varepsilon_ {i _ {1}} + \dots + \varepsilon_ {i _ {k}})) (h) (e _ {- i _ {1}} \wedge \dots \wedge e _ {- i _ {k}}). \end{array} \tag {5}
\]

这表明 \(\rho\) 的最高权是 \(\varpi_{l}\)。我们称 \(\rho\) 为 g 的旋量表示。注意，其权全为单权（此外，\(\varpi_{l}\) 是极小权）。

\begin{enumerate}
\def\labelenumi{(\Alph{enumi})}
\setcounter{enumi}{21}
\tightlist
\item
有 \(w_0 = -1\)，所以 \(\mathfrak{g}\) 的每个有限维单表示都是正交型或辛型。由第 VI 章的结果，\(\varpi_r\) 关于 \((\alpha_1, \ldots, \alpha_l)\) 的坐标之和在 \(1 \leq r \leq l - 1\) 时为整数；因此，表示 \(\bigwedge^r \sigma\) 是正交型的。此外，它保持 \(\Psi_{(r)}\)（\(\Psi\) 在 \(\bigwedge^r V\) 上的延拓）不变。
\end{enumerate}

对于旋量表示，\(\varpi_{l}\) 关于 \((\alpha_{1},\ldots,\alpha_{l})\) 的坐标之和为 \(\frac{1}{2}(1+\cdots+l)=\frac{l(l+1)}{4}\)（同上引）。因此，当 \(l\equiv0\) 或 -1 (mod. 4) 时它是正交型的，当 \(l\equiv1\) 或 2 (mod. 4) 时它是辛型的。事实上，考虑 \(\Phi\) 上的双线性型 \(N=\bigwedge F'\)，定义如下：若 \(x\in\bigwedge^{p}F'\) 且 \(y\in\bigwedge^{q}F'\)，则当 \(\Phi(x,y)=0\) 时置 \(p+q\neq l\)，并且

\[
x \wedge y = (- 1) ^ {\frac {p (p + 1)}{2}} \varPhi (x, y) e _ {- 1} \wedge \dots \wedge e _ {- r}
\]

当 \(p + q = l\) 时，容易验证 \(\Phi\) 非退化；当 \(l \equiv 0, -1 \pmod{4}\) 时，\(l \equiv 1, 2 \pmod{4}\) 也是非退化的。

\[
f (X _ {\varepsilon_ {i}}) = e _ {0} e _ {i}, \quad f (X _ {- \varepsilon_ {i}}) = - e _ {0} e _ {- i}
\]

对 \(1 \leq i \leq l\)，且

\[
f (X _ {\varepsilon_ {i} - \varepsilon_ {j}}) = \frac {1}{2} (e _ {i} e _ {- j} - e _ {- j} e _ {i}) = e _ {i} e _ {- j} = e _ {0} e _ {i} e _ {0} e _ {- j}
\]

对于 \(1 \leq i < j \leq l\)，同理

\[
f (X _ {\varepsilon_ {j} - \varepsilon_ {i}}) = - e _ {0} e _ {j} e _ {0} e _ {- i}, f (X _ {\varepsilon_ {i} + \varepsilon_ {j}}) = e _ {0} e _ {i} e _ {0} e _ {j}, f (X _ {- \varepsilon_ {i} - \varepsilon_ {j}}) = e _ {0} e _ {- i} e _ {0} e _ {- j};
\]

从而

\[
\mu^ {- 1} \circ f (X _ {\varepsilon_ {i}}) = e _ {i}, \quad \mu^ {- 1} f (X _ {- \varepsilon_ {i}}) = - e _ {- i}
\]

且

\[
\mu^ {- 1} \circ f (X _ {\pm \varepsilon_ {i} \pm \varepsilon_ {j}}) = c e _ {\pm i} e _ {\pm j} \quad \text {   for   } 1 \leq i, j \leq l, i \neq j \text {   with   } c \in \{1, - 1 \}.
\]

现在很容易验证，\(\Phi\) 实际上在 g 下不变（参见~习题18）。

\begin{enumerate}
\def\labelenumi{(\Roman{enumi})}
\setcounter{enumi}{5}
\tightlist
\item
对于 \(x \in \mathfrak{g}\)，\(\sigma(x)\) 的特征多项式具有形式
\end{enumerate}

\[
\mathrm{T} ^ {2 l + 1} + f _ {1} (x) \mathrm{T} ^ {2 l} + f _ {2} (x) \mathrm{T} ^ {2 l - 1} + \dots + f _ {2 l + 1} (x)
\]

其中 \(f_{1},\ldots ,f_{2l + 1}\) 是 \(\mathfrak{g}\) 上的不变多项式函数。

若 \(x = \xi_{1}H_{1} + \cdots + \xi_{l}H_{l} \in h\)，则 \(f_{i}(x)\) 在相差一个符号的意义下是 \(\xi_{1}, \ldots, \xi_{l}, -\xi_{1}, \ldots, -\xi_{l}\) 的初等对称函数；这些对称函数在奇数次数时为零，并且

\[
\mathrm{T} ^ {2 l + 1} + f _ {2} (x) \mathrm{T} ^ {2 l - 1} + f _ {4} (x) \mathrm{T} ^ {2 l - 3} + \dots + f _ {2 l} (x) \mathrm{T} = \mathrm{T} (\mathrm{T} ^ {2} - \xi_ {1} ^ {2}) \dots (\mathrm{T} ^ {2} - \xi_ {l} ^ {2})
\]

从而 \(f_{2},\ldots,f_{2l}\) 在相差一个符号的意义下是 \(\xi_{1}^{2},\ldots,\xi_{l}^{2}\) 的初等对称函数，它们是 \(\mathbf{S}(\mathfrak{h}^{*})^{\mathrm{W}}\) 的代数无关生成元（第 VI 章，§4，第5.IX节）。根据 §8，第3节，定理1 (i)，可知 \(f_{1}=f_{3}=f_{5}=\cdots=0\)，并且 \((f_{2},f_{4},\ldots,f_{2l})\) 是生成 g 上不变多项式函数代数的代数自由族。

\begin{enumerate}
\def\labelenumi{(\Roman{enumi})}
\setcounter{enumi}{6}
\tightlist
\item
由于邓金图的唯一自同构是恒等自同构，有 \(\mathrm{Aut}(\mathfrak{g}) = \mathrm{Aut}_0(\mathfrak{g})\)。
\end{enumerate}

令 \(\Sigma\) 为关于 \(\Psi\) 的 V 的相似变换群。对所有 \(g \in \Sigma\)，令 \(\varphi(g)\) 为自同构 \(x \mapsto gxg^{-1}\)，它作用于 \(\mathfrak{g}\)。于是 \(\varphi\) 是从 \(\Sigma\) 到 \(\operatorname{Aut}(\mathfrak{g})\) 的同态。我们证明它是满射。令 \(\alpha \in \operatorname{Aut}(\mathfrak{g}) = \operatorname{Aut}_0(\mathfrak{g})\)。由 §7，第1节，命题2，存在 \(s \in \mathbf{GL}(V)\)，使得 \(\alpha(x) = sxs^{-1}\) 对所有 \(x \in \mathfrak{g}\) 成立。于是 \(s\) 将 \(\Psi\) 变为双线性型 \(\Psi'\)，它在 \(\mathfrak{g}\) 下不变，从而与 \(\Psi\) 成比例；这证明 \(s \in \Sigma\)。

由于 g 的恒等表示不可约，其交换子仅为标量（\(\S6\)，第1节，命题1），所以 \(\varphi\) 的核为 \(k^{*}\)。因此，群 \(\operatorname{Aut}(\mathfrak{g}) = \operatorname{Aut}_{0}(\mathfrak{g})\) 可与 \(\Sigma/k^{*}\) 识别。但是，由《代数》第 IX 章，\(\S6\)，第5节，\(\Sigma\) 是 \(k^{*}\) 与 \(\mathbf{SO}(\Psi)\) 的乘积；所以 \(\operatorname{Aut}(\mathfrak{g}) = \operatorname{Aut}_{0}(\mathfrak{g})\) 可与 \(\mathbf{SO}(\Psi)\) 识别。

令 \(\mathbf{O}_0^+(\varPsi)\) 为 \(\varPsi\) 的既约正交群（《代数》第 IX 章，§9，第5节）。由于 \(\mathbf{SO}(\varPsi)/\mathbf{O}_0^+(\varPsi)\) 是交换群（同上引），群 \(\mathrm{Aut}_e(\mathfrak{g})\) 包含于 \(\mathbf{O}_0^+(\varPsi)\)（§11，第2节，命题3）；事实上它等于该群（习题7）。

\begin{enumerate}
\def\labelenumi{(\Roman{enumi})}
\setcounter{enumi}{7}
\tightlist
\item
\(\Phi_{R}\) 上的典则双线性型 \(h^{*}\) 为
\end{enumerate}

\[
\varPhi_ {\mathrm{R}} (\xi_ {1} \varepsilon_ {1} + \dots + \xi_ {l} \varepsilon_ {l}, \xi_ {1} ^ {\prime} \varepsilon_ {1} + \dots + \xi_ {l} ^ {\prime} \varepsilon_ {l}) = \frac {1}{4 l - 2} (\xi_ {1} \xi_ {1} ^ {\prime} + \dots + \xi_ {l} \xi_ {l} ^ {\prime})
\]

（第 VI 章，§4，第5.V节）。由 \(\mathfrak{h}\) 定义的从 \(\mathfrak{h}^*\) 到 \(\varPhi_{\mathrm{R}}\) 的同构将 \(H_{i}\) 映为 \((4l - 2).\varepsilon_{i}\)。因此，\(\varPhi_{\mathrm{R}}\) 的逆型，即 Killing 型在 \(\mathfrak{h}\) 上的限制，为

\[
\varPhi (\xi_ {1} H _ {1} + \dots + \xi_ {l} H _ {l}, \xi_ {1} ^ {\prime} H _ {1} + \dots + \xi_ {l} ^ {\prime} H _ {l}) = (4 l - 2) (\xi_ {1} \xi_ {1} ^ {\prime} + \dots + \xi_ {l} \xi_ {l} ^ {\prime}).
\]

\begin{enumerate}
\def\labelenumi{(\Roman{enumi})}
\setcounter{enumi}{8}
\tightlist
\item
回顾由公式（3）定义的 \(X_{\alpha}\)（\(\alpha \in R\)）。容易验证，对于 \([X_{\alpha}, X_{-\alpha}] = -H_{\alpha}\)，有 \(\alpha \in R\)。另一方面，令 M 为矩阵 \(I + E_{0,0}\)；由于 \(M = S^{t}M^{-1}S\)，映射
\end{enumerate}

\[
\theta : g \mapsto - M ^ {- 1 t} g M
\]

是 \(\mathfrak{g}\) 的自同构，且 \(\theta(X_{\alpha}) = X_{-\alpha}\) 对所有 \(\alpha \in \mathrm{R}\) 成立。因此，\((X_{\alpha})\) 是 \((\mathfrak{g},\mathfrak{h})\) 中的谢瓦莱系。

假设 \(k = \mathbf{Q}\)。嘉当子代数 \(\mathfrak{h}\) 有两个许可格：格 \(\mathrm{Q}(\mathrm{R}^{\vee})\) 由 \(H_{\alpha}\) 生成，格 \(\mathrm{P}(\mathrm{R}^{\vee})\) 由 \(H_{i}\) 生成，且由 \(\mathfrak{h}\) 中的整数项对角矩阵组成。因此，\(\mathfrak{o}_S(2l + 1, \mathbf{Z})\)（即 \(\mathfrak{g}\) 中具有整数项的矩阵集合）是谢瓦莱序 \(\mathrm{P}(\mathrm{R}^{\vee}) + \sum \mathbf{Z}.X_{\alpha}\)（属于 \(\mathfrak{g}\)）。由于 \((X_{\pm \varepsilon_i})^2 = 2E_{\pm i,\mp i}, (X_{\pm \varepsilon_i})^3 = 0\)，且 \((X_{\pm \varepsilon_i\pm \varepsilon_j})^2 = 0\)，可知格 \(\mathcal{V}\) 由 Witt 基 \((e_i)_{-l\leq i\leq l}\) 生成，是 \(\mathfrak{o}_S(2l + 1,\mathbf{Z})\) 中关于 \(V\) 的容许格。\(\bigwedge^r\mathcal{V}\) 在 \(\bigwedge^r V\) 中也如此。

现在考虑 g 在 \(\rho\) 上的旋量表示 \(N = \bigwedge F'\)。由于其各权不将 \(\mathrm{P}(\mathrm{R}^{\vee})\) 映入 Z，它对 \(\mathfrak{o}_{S}(2l + 1, \mathbf{Z})\) 没有容许格。另一方面，由 N 的典则基 \((e_{-i_{1}} \wedge \cdots \wedge e_{-i_{k}})\)（\(1 \leq i_{1} < \cdots < i_{k} \leq l\)）生成的格 N 是谢瓦莱序 \(\mathcal{G} = \mathrm{Q}(\mathrm{R}^{\vee}) + \sum_{\alpha \in \mathrm{R}} \mathbf{Z}. X_{\alpha}\) 的容许格。事实上，显然 N 在以 \(e_{-i}\) 作外积以及以 \(e_{i}\) 作内积下稳定（\(1 \leq i \leq l\)）。(V) 中的公式于是表明 N 在 \(\rho(\mathcal{G})\) 下稳定。此外，由于对所有 \(\rho(X_{\alpha})^{2} = 0\) 都有 \(\alpha \in R\)，可知 N 是容许格。

